Para el hiperboloide de dos hojas $\dfrac{z^2}{4} - \dfrac{x^2}{9} - \dfrac{y^2}{9} = 1$:
\begin{enumerate}
    \item Determinar los valores de $z$ para los cuales existe la superficie. ¿Cuál es la región vacía?
    \item Encontrar los vértices del hiperboloide (los puntos de las hojas más cercanos al origen).
    \item Encontrar la traza con el plano $z = 4$ e identificar su tipo y semiejes.
    \item Encontrar las trazas con $y = 0$ y $x = 0$. ¿Qué hipérbolas se obtienen?
    \item Comparar las ecuaciones del hiperboloide de una hoja y del de dos hojas del ejercicio anterior. ¿Qué diferencia hay en los signos? ¿Cómo afecta esto a la topología (número de componentes conexas)?
\end{enumerate}